% Методичка требует краткие теоретические сведения, поэтому раздел включён.
\section{Теоретическая часть}
Совершенная дизъюнктивная нормальная форма — дизъюнкция конъюнкций, по одной на каждый
набор входов, на котором функция равна 1. В каждую конъюнкцию входят все переменные: без
отрицания, если на наборе переменная равна 1, и с отрицанием, если 0. СДНФ однозначно
задаёт функцию, но обычно избыточна.

Карта Карно — таблица значений функции, у которой строки и столбцы подписаны кодом Грея:
соседние клетки, в том числе через край карты, отличаются значением одной переменной. Две
соседние единицы склеиваются в одну конъюнкцию без этой переменной. Контур из $2^k$ единиц
даёт конъюнкцию, в которой на $k$ переменных меньше. Минимальная ДНФ получается из
наименьшего числа наибольших контуров, покрывающих все единицы.
